\subsection{The prime spectrum of a ring}

Let A be a ring and X the set of prime ideals of A. For every subset M of A, we denote by V(M) the set of prime ideals of A containing M; clearly, if a is the ideal of A generated by M, V(M) = V(a); if M consists of a single element f, we write V(f) instead of V(\{f\}) and we have V(f) = V(Af). The mapping M \(\mapsto\) V(M) is decreasing with respect to inclusion in A and X. Moreover, the following formulae hold:

\begin{enumerate}
\def\labelenumi{(\arabic{enumi})}
\tightlist
\item
\end{enumerate}

\[
\mathbf {V} (0) = \mathbf {X}, \quad \mathbf {V} (1) = \varnothing ;\tag{1}
\]

\[
\mathrm{V} \left(\bigcup_ {i \in I} \mathrm{M} _ {i}\right) = \mathrm{V} \left(\sum_ {i \in I} \mathrm{M} _ {i}\right) = \bigcap_ {i \in I} \mathrm{V} (\mathrm{M} _ {i})\tag{2}
\]

for every family \((\mathbf{M}_i)_{i\in I}\) of subsets of A;

\[
\mathbf {V} (\mathfrak {a} \cap \mathfrak {a} ^ {\prime}) = \mathbf {V} (\mathfrak {a a} ^ {\prime}) = \mathbf {V} (\mathfrak {a}) \cup \mathbf {V} (\mathfrak {a} ^ {\prime})\tag{3}
\]

for every pair of ideals \(\mathfrak{a},\mathfrak{a}'\) in A. Formulae (1) and (2) are obvious; on the other hand, formula (3) means that, for a prime ideal \(\mathfrak{p}\) of A to contain one of the ideals \(\mathfrak{a}\) or \(\mathfrak{a}'\) , it is necessary and sufficient that it contain \(\mathfrak{a}\mathfrak{a}'\) or that it contain \(\mathfrak{a} \cap \mathfrak{a}'\) ; then it is a consequence of § 1, no. 1, Proposition 1. The second formula (1) has the following converse: if \(\mathfrak{a}\) is an ideal of \(\mathbf{A}\) such that \(V(\mathfrak{a}) = \varnothing\) , then \(\mathfrak{a} = \mathbf{A}\) , for there is no maximal ideal of \(\mathbf{A}\) containing \(\mathfrak{a}\) . Finally, if \(\mathfrak{a}\) is an ideal of \(\mathbf{A}\) and \(r(\mathfrak{a})\) is its radical (§ 2, no. 6, Definition 4), then

\[
\mathbf {V} (\mathfrak {a}) = \mathbf {V} (\mathfrak {r} (\mathfrak {a}))\tag{4}
\]

as follows from § 2, no. 6, Corollary 1 to Proposition 13.

Formulae (1) to (3) show that the subsets V(M) of X satisfy the closed set axioms of a topology (General Topology, Chapter I, § 1, no. 4).

DEFINITION 4. Let A be a ring. The set X of prime ideals of A, with the topology whose closed sets are the sets V(M), where M runs through P(A), is called the prime spectrum of A and denoted by Spec(A). The topology thus defined is called the spectral or Zariski topology on X.

Clearly the relation \(\operatorname{Spec}(\mathbf{A}) = \varnothing\) is equivalent to \(\mathbf{A} = \{0\}\) .

Let X be the prime spectrum of a ring A; for all \(f \in A\) , let us denote by \(X_{f}\) the set of prime ideals of A not containing f; then \(\mathbf{X}_{f} = \mathbf{X} - \mathbf{V}(f)\) and \(X_{f}\) is therefore an open set. By (2), every closed subset of X is an intersection of closed sets of the form \(\mathrm{V}(f)\) and hence the \(X_{f}\) form a base of the spectral topology on X. Moreover, it follows immediately from the definitions that

\[
\mathbf {X} _ {0} = \varnothing , \quad \mathbf {X} _ {1} = \mathbf {X},\tag{5}
\]

and more generally \(X_{f} = X\) for every invertible element f of A;

\[
\mathrm{X} _ {f g} = \mathrm{X} _ {f} \cap \mathrm{X} _ {g} \quad \text {   for   all   } f, g \text {   in   A.   }\tag{6}
\]

For every subset Y of X, let us denote by \(\Im(Y)\) the intersection of the prime ideals of A which belong to Y. Clearly \(\Im(Y)\) is an ideal of A and the mapping \(Y \mapsto \Im(Y)\) is decreasing with respect to inclusion in X and A. Clearly the relations

\begin{enumerate}
\def\labelenumi{(\arabic{enumi})}
\setcounter{enumi}{6}
\tightlist
\item
\end{enumerate}

\[
\Im (\varnothing) = \mathbf {A}\tag{7}
\]

\[
\Im \left(\bigcup_ {\lambda \in L} Y _ {\lambda}\right) = \bigcap_ {\lambda \in L} \Im (Y _ {\lambda})\tag{8}
\]

hold for every family \((\mathbf{Y}_{\lambda})_{\lambda \in \mathbf{L}}\) of subsets of \(\mathbf{X}\) . Moreover:

PROPOSITION 11. Let A be a ring, a an ideal of A and Y a subset of X = Spec(A).

\begin{enumerate}
\def\labelenumi{(\roman{enumi})}
\item
  \(\mathbf{V}(\mathfrak{a})\) is closed in \(\mathbf{X}\) and \(\Im (\mathbf{Y})\) is an ideal of A which is equal to its radical.
\item
  \(\Im(\mathrm{V}(\mathfrak{a}))\) is the radical of \(\mathfrak{a}\) and \(\mathrm{V}(\Im(\mathrm{Y}))\) is the closure of \(\mathrm{Y}\) in \(\mathrm{X}\) .
\item
  The mappings \(\Im\) and \(\mathbf{V}\) define decreasing bijections, one of which is the inverse of the other, between the set of closed subsets of \(\mathbf{X}\) and the set of ideals of \(\mathbf{A}\) which are equal to their radicals.
\end{enumerate}

Assertion (i) and the first assertion of (ii) follow from the definitions and § 2, no. 6, Corollary 1 to Proposition 13. If a closed set V(M) (for some M ⊂ A) contains Y, then M ⊂ p for every prime ideal p ∈ Y, whence M ⊂ \(\mathfrak{I}\)(Y) and consequently V(M) ⊃ V(\(\mathfrak{I}\)(Y)); as Y ⊂ V(\(\mathfrak{I}\)(Y)), V(\(\mathfrak{I}\)(Y)) is the smallest closed set of X containing Y, which completes the proof of (ii). Finally, it follows from (ii) that, if a is an ideal equal to its radical, then \(\mathfrak{I}\)(V(a)) = a and that, if Y is closed in X, then V(\(\mathfrak{I}\)(Y)) = Y; this proves (iii).

It follows immediately from Proposition 11 that, if M is any subset of A and Y is any subset of X, then V(M) = V( \(\Im(V(M))\) ) and \(\Im(Y) = \Im(V(\Im(Y)))\) .

COROLLARY 1. For every family \((\mathbf{Y}_{\lambda})_{\lambda \in L}\) of closed subsets of \(\mathbf{X}\) , \(\Im\left(\bigcap_{\lambda \in L} \mathbf{Y}_{\lambda}\right)\) is the radical of the sum of the ideals \(\Im(\mathbf{Y}_{\lambda})\) .

It follows from Proposition 11 (iii) that \(\Im\left(\bigcap_{\lambda\in L}Y_{\lambda}\right)\) is the smallest ideal which is equal to its radical and contains all the \(\Im(Y_{\lambda})\) ; this ideal then contains \(\sum_{\lambda\in L}\Im(Y_{\lambda})\) and therefore also the radical of \(\sum_{\lambda\in L}\Im(Y_{\lambda})\) ( \(\S\) 2, no. 6, Corollary 2 to Proposition 13), whence the corollary.

COROLLARY 2. Let \(\mathfrak{r}(\mathfrak{a})\) denote the radical of an ideal \(\mathfrak{a}\) of \(\mathbf{A}\) ; if \(\mathfrak{a}\) and \(\mathfrak{b}\) are two ideals of \(\mathbf{A}\) , the relation \(\mathrm{V}(\mathfrak{a}) \subset \mathrm{V}(\mathfrak{b})\) is equivalent to \(\mathfrak{b} \subset \mathfrak{r}(\mathfrak{a})\) and \(\mathfrak{r}(\mathfrak{b}) \subset \mathfrak{r}(\mathfrak{a})\) .

It is immediate that the relations \(\mathfrak{b} \subset \mathfrak{r}(\mathfrak{a})\) and \(r(\mathfrak{b}) \subset \mathfrak{r}(\mathfrak{a})\) are equivalent and, as \(\mathrm{V}(\mathfrak{a}) = \mathrm{V}(\mathfrak{r}(\mathfrak{a}))\) , the corollary follows immediately from Proposition 11, (iii).

COROLLARY 3. Let \((f_{\lambda})_{\lambda \in L}\) be a family of elements of A. For an element \(g \in \mathbf{A}\) to satisfy \(\mathbf{X}_g \subset \bigcup_{\lambda \in L} \mathbf{X}_{f_\lambda}\) , it is necessary and sufficient that there exist an integer \(n > 0\) such that \(g^n\) belongs to the ideal generated by the \(f_\lambda\) .

The relation \(X_{g} \subset \bigcup_{\lambda \in L} X_{f_{\lambda}}\) is equivalent to \(\mathrm{V}(g) \supset \bigcap_{\lambda \in L} \mathrm{V}(f_{\lambda})\) and it is sufficient to apply Corollary 2.

COROLLARY 4. For two elements \(f, g\) of \(A\) to satisfy \(X_f = X_g\) , it is necessary and sufficient that there exist two integers \(m > 0, n > 0\) such that \(f^m \in Ag\) and \(g^n \in Af\) .

* COROLLARY 5. For \(f \in \mathbf{A}\) to satisfy \(\mathbf{X}_f = \varnothing\) , it is necessary and sufficient that \(f\) be nilpotent.

This follows immediately from Corollary 4.

COROLLARY 6. The closure of a set consisting of a point \(\mathfrak{p} \in \mathbf{X} = \operatorname{Spec}(\mathbf{A})\) is the set \(\mathbf{V}(\mathfrak{p})\) of prime ideals containing \(\mathfrak{p}\) . For the set \(\{\mathfrak{p}\}\) to be closed in \(\mathbf{X}\) (or, as we shall also say by an abuse of language, for \(\mathfrak{p}\) to be a closed point of \(\mathbf{X}\) ), it is necessary and sufficient that \(\mathfrak{p}\) be maximal.

COROLLARY 7. If A is a Noetherian ring, X = Spec(A) is a Noetherian space.

PROPOSITION 12. For all \(f \in \mathbf{A}\) , the open set \(\mathbf{X}_f\) in \(\mathbf{X} = \operatorname{Spec}(\mathbf{A})\) is quasi-compact; in particular, the space \(\mathbf{X}\) is quasi-compact.

As the \(X_{g}\) form a base of the topology, it is sufficient to prove that, if \((g_{\lambda})_{\lambda\in L}\) is a family of elements of A such that \(X_{f}\subset\bigcup_{\lambda\in L}X_{g_{\lambda}}\) , then there exists a finite subfamily \((g_{\lambda})_{\lambda\in H}\) such that \(X_{f}\subset\bigcup_{\lambda\in H}X_{g_{\lambda}}\) . But the relation \(X_{f}\subset\bigcup_{\lambda\in L}X_{g_{\lambda}}\) means that there exists an integer \(n > 0\) and a finite subfamily \((g_{\lambda})_{\lambda\in L}\) such that \(f^{n}\) belongs to the ideal generated by that subfamily (Corollary 3 to Proposition 11); whence the proposition.

PROPOSITION 13. Let A, A' be two rings, X = Spec(A), X' = Spec(A') and h a homomorphism from A to A'; the mapping \(^a h\) : \(p' \mapsto \overline{h}^{-1}(p')\) from X' to X is continuous.

For \(\mathbf{M} \subset \mathbf{A}\) , the set \((^a h)^{-1}(\mathbf{V}(\mathbf{M}))\) is the set of prime ideals \(\mathfrak{p}'\) of \(\mathbf{A}'\) such that \(\mathbf{M} \subset \overline{h}^{-1}(\mathfrak{p}')\) , which is equivalent to \(h(\mathbf{M}) \subset \mathfrak{p}'\) ; this set is then equal to \(\mathbf{V}(h(\mathbf{M}))\) and is therefore closed.

We call \(^{a}h\) the mapping associated with the homomorphism h.

Remark. If \(h\) is surjective and \(\mathfrak{a}\) is its kernel, it follows from the definition of the spectral topology that \(^a h\) : is a homeomorphism of \(X'\) onto the closed subspace \(V(\mathfrak{a})\) of \(X\) ; for a prime ideal \(\mathfrak{p}'\) of \(A'\) to contain an ideal \(\mathfrak{b}'\) of \(A'\) , it is necessary and sufficient that \(\overline{h}^{-1}(\mathfrak{p}')\) contain \(\overline{h}^{-1}(\mathfrak{b}')\) ; we see first that \(^a h\) is injective by taking \(\mathfrak{b}'\) to be prime; moreover, for every ideal \(\mathfrak{b}'\) of \(A'\) , the image under \(^a h\) of \(V(\mathfrak{b}')\) is \(V(\overline{h}^{-1}(\mathfrak{b}'))\) , whence our assertion, the ideals of the form \(\overline{h}^{-1}(\mathfrak{b}')\) all being ideals of \(A\) which contain \(\mathfrak{a}\) .

COROLLARY. Let S be a multiplicative subset of A, \(A' = S^{-1}A\) and h the canonical homomorphism \(i_{A}^{S}\) ; then \(^{a}h\) is a homeomorphism of \(X' = \operatorname{Spec}(A')\) onto the subspace of \(X = \operatorname{Spec}(A)\) consisting of the prime ideals of A which do not meet S.

Let \(f' = f / s\) , where \(f \in \mathbf{A}\) , \(s \in \mathbf{S}\) ; then \(\mathbf{X}_{f'}' = \mathbf{X}_{f/1}'\) since \(s / 1\) is invertible in \(\mathbf{A}'\) . We know already that \(^a h\) is injective and that, for all \(p' \in \mathbf{X}'\) , the relations \(f / 1 \in p'\) and \(f \in \bar{h}^{-1}(p') = ^a h(p')\) are equivalent and hence the conditions \(p' \in \mathbf{X}_{f/1}'\) and \(^a h(p') \in \mathbf{X}_f\) are equivalent; this shows that \(^a h(\mathbf{X}_{f'})\) is equal to \(\mathbf{X}_f \cap ^a h(\mathbf{X}')\) , whence the first assertion, since the \(\mathbf{X}_f\) (resp. \(\mathbf{X}_{f'}'\) ) form a base of the topology of \(\mathbf{X}\) (resp. \(\mathbf{X}'\) ). The second assertion follows from § 2, no. 5, Proposition 11 (ii).

PROPOSITION 14. Let A be a ring. For a subset Y of X = Spec(A) to be irreducible, it is necessary and sufficient that the ideal \(\Im(Y)\) be prime.

Writing \(\mathfrak{p} = \Im(\mathbf{Y})\) , we note that, for an element \(f \in \mathbf{A}\) , the relation \(f \in \mathfrak{p}\) is equivalent to \(\mathbf{Y} \subset \mathbf{V}(f)\) . Suppose that \(\mathbf{Y}\) is irreducible and let \(f, g\) be elements of \(\mathbf{A}\) such that \(fg \in \mathfrak{p}\) . Then

\[
\mathbf {Y} \subset \mathbf {V} (f g) = \mathbf {V} (f) \cup \mathbf {V} (g);
\]

as Y is irreducible and V(f) and V(g) are closed, Y ⊂ V(f) or Y ⊂ V(g), hence f ∈ p or g ∈ p, which proves that p is prime.

Suppose now that p is prime; then \(\overline{\mathbf{Y}} = \mathbf{V}(\mathfrak{p})\) (Proposition 11 (ii)) and, as p is prime, \(\mathfrak{p} = \Im(\{\mathfrak{p}\})\) , whence \(\overline{\mathbf{Y}} = \mathbf{V}(\Im(\{\mathfrak{p}\})) = \{\mathfrak{p}\}\) (Proposition 11 (ii)). As a set consisting of a single point is irreducible, Y is irreducible (no. 1, Proposition 2).

COROLLARY 1. For a ring A to be such that \(\mathbf{X} = \operatorname{Spec}(\mathbf{A})\) is irreducible, it is necessary and sufficient that the quotient of A by its nilradical \(\mathfrak{N}\) be an integral domain.

By Proposition 11 (i), \(\Im(\mathbf{X})\) is the radical of the ideal (0), that is \(\mathfrak{N}\) .

COROLLARY 2. The mapping \(\mathfrak{p} \mapsto \mathrm{V}(\mathfrak{p})\) is a bijection of \(\mathrm{X} = \operatorname{Spec}(\mathrm{A})\) onto the set of irreducible closed subsets of \(\mathrm{X}\) ; in particular the irreducible components of a closed subset \(\mathrm{Y}\) of \(\mathrm{X}\) are the sets \(\mathrm{V}(\mathfrak{p})\) , where \(\mathfrak{p}\) runs through the set of minimal elements of the set of prime ideals of \(\mathrm{A}\) which contain \(\Im(\mathrm{Y})\) .

As \(\Im (\mathbf{V}(\mathfrak{p})) = \mathfrak{p}\) for every prime ideal \(\mathfrak{p}\) of \(\mathbf{A}\) and \(\mathbf{Y} = \mathbf{V}(\Im (\mathbf{Y}))\) for every closed subset \(\mathbf{Y}\) of \(\mathbf{X}\) , the first assertion follows from Proposition 14; on the other hand, for \(\mathbf{Y} \supset \mathbf{V}(\mathfrak{p})\) , it is necessary and sufficient that

\[
\mathfrak {p} = \Im (\mathbf {V} (\mathfrak {p})) \supset \Im (\mathbf {Y})
\]

(Proposition 11), whence the second assertion.

COROLLARY 3. The set of minimal prime ideals of a Noetherian ring A is finite.

X = Spec(A) has only a finite number of irreducible components (Corollary 7 to Proposition 11 and no. 2, Proposition 10) and the corollary follows from Corollary 2 above.

PROPOSITION 15. Let A be a ring, I a finite set and E the set of orthogonal families \((e_i)_{i \in I}\) of idempotents \(e_i \neq 0\) of A such that \(\sum_{i \in I} e_i = 1\) . For all \((e_i)_{i \in I} \in E\) , we set \(\overline{\omega}((e_i)_{i \in I}) = (V(A(1 - e_i)))_{i \in I}, \sigma((e_i)_{i \in I}) = (Ae_i)_{i \in I}\) . Then \(\overline{\omega}\) is a bijection of E onto the set P of partitions \((U_i)_{i \in I}\) of \(X = \text{Spec}(A)\) into open sets and \(\sigma\) is a bijection of E onto the set S of families \((a_i)_{i \in I}\) of ideals \(\neq 0\) of A such that A is the direct sum of the \(a_i\) .

Let \((e_i)_{i\in I}\) be an element of \(E\) and set \(\mathbf{Y}_i = \mathbf{V}(\mathbf{A}(1 - e_i))\) ; if \(i\neq j\) , then \(1 = 1 - e_i + e_i(1 - e_j)\in A(1 - e_i) + A(1 - e_j)\) , whence \(\mathbf{Y}_i\cap \mathbf{Y}_j = \varnothing\) (formulae (1) and (2)). On the other hand,

\[
\bigcup_ {i \in I} Y _ {i} = V \left(\prod_ {i \in I} A (1 - e _ {i})\right) \quad (\text { formula   (3) });
\]

by hypothesis \(\prod_{i\in I}(1 - e_i) = 1 - \sum_{i\in I}e_i = 0\) , whence \(\bigcup_{i\in I}\mathbf{Y}_i = \mathbf{X}\) (formula (1)). As the \(\mathbf{Y}_i\) are closed, they are also open, whence \(\overline{\omega} (\mathrm{E})\subset \mathrm{P}\) . Also, obviously \(\mathbf{A} = \sum_{i\in I}\mathbf{A}e_i\) ; if \(0 = \sum_{i\in I}a_{i}e_{i}\) where \(a_{i}\in \mathbf{A}\) , we obtain, by multiplying by \(e_i\) , \(0 = a_i e_i^2 = a_i e_i\) for all \(i\) ; this proves that \(\sigma (\mathrm{E})\subset \mathrm{S}\) .

LEMMA 2. If \(e, f\) are two idempotents of \(A\) such that \(Ae\) and \(Af\) have the same radical then \(e = f\) .

There exists by hypothesis integers \(m \geqslant 0\) , \(n \geqslant 0\) such that \(e = e^{m} \in Af\) and \(f = f^{n} \in Ae\) ; let x, y be elements of A such that e = xf, f = ye; then ef = \(xf^{2} = xf = e\) and similarly ef = \(ye^{2} = ye = f\) , whence e = f.

Lemma 2 and Corollary 2 to Proposition 11 show that the mappings \(\overline{\omega}\) and \(\sigma\) are injective.

Let us show that \(\sigma\) is surjective. If \((a_{i})_{i\in I}\) is an element of \(S\) , there are elements \(e_i\in a_i\) such that \(1 = \sum_{i\in I}e_i\) ; if \(i\neq j\) , then \(e_i e_j\in a_i\cap a_j = \{0\}\) , whence

\[
e _ {i} = \sum_ {j \in I} e _ {i} e _ {j} = e _ {i} ^ {2};
\]

finally, \(\mathbf{A}e_{i}\subset \mathfrak{a}_{i}\) for all \(i\in \mathbf{I}\) and \(\sum_{i\in \mathbf{I}}\mathrm{A}e_i = \mathrm{A},\) whence \(\mathrm{A}e_{i} = \mathfrak{a}_{i}\) .

It remains to prove that \(\overline{\omega}\) is surjective. Let \((\mathbf{U}_i)_{i\in I}\) be an element of \(P\) and set \(Z_{i} = \complement U_{i} = \bigcup_{j\neq i}U_{j}\) ; as \(U_{i}\) and \(Z_{i}\) are closed, there exist ideals \(\mathfrak{a}_i,\mathfrak{b}_i\) of \(A\) such that \(\mathbf{U}_i = \mathbf{V}(\mathfrak{a}_i)\) , \(Z_{i} = \mathbf{V}(\mathfrak{b}_{i})\) . We now show that it is possible to suppose further that \(\mathfrak{a}_i\cap \mathfrak{b}_i = 0\) . Now \(\mathbf{U}_i\cap Z_i = \varnothing\) , whence \(\mathfrak{a}_i + \mathfrak{b}_i = A\) ; let \(a_i\in \mathfrak{a}_i\) , \(b_{i}\in \mathfrak{b}_{i}\) be such that \(a_{i} + b_{i} = 1\) . Then \(\mathbf{X} = \mathbf{U}_i\cup Z_i = \mathbf{V}(\mathfrak{a}_i\mathfrak{b}_i)\) (formula (3)); every element of \(\mathfrak{a}_i\mathfrak{b}_i\) is therefore nilpotent (Corollary 2 to Proposition 11); let \(p\) be an integer such that \(a_i^p b_i^p = 0\) . Then \(\mathbf{U}_i\subset \mathbf{V}(\mathbf{A}a_i) = \mathbf{V}(\mathbf{A}a_i^p)\) ,

\[
\mathbf {Z} _ {i} \subset \mathbf {V} (\mathbf {A} b _ {i}) = \mathbf {V} (\mathbf {A} b _ {i} ^ {p})
\]

and \(\mathrm{V}(\mathbf{A}a_i)\cap \mathrm{V}(\mathbf{A}b_i) = \mathrm{V}(\mathbf{A}a_i + \mathbf{A}b_i) = \varnothing\) , hence \(\mathbf{U}_i = \mathrm{V}(\mathbf{A}a_i^p)\) and \(\mathbf{Z}_i = \mathrm{V}(\mathbf{A}b_i^p)\) , which establishes our assertion by replacing \(a_i\) by \(\mathbf{A}a_i^p\) and \(b_i\) by \(\mathbf{A}b_i^p\) . The ideals \(a_i\) and \(b_i\) thus chosen, it follows from the fact that \(\sigma\) is bijective that there exist two idempotents \(f_i\in a_i\) , \(e_i\in b_i\) such that \(1 = e_i + f_i\) , \(e_i f_i = 0\) , \(a_i = Af_i\) , \(b_i = Ae_i\) . If \(i\neq j\) , then \(\mathbf{X} = \mathbf{Z}_i\cup \mathbf{Z}_j = \mathrm{V}(\mathbf{A}e_i e_j)\) , and as \(e_i e_j\) is idempotent, Lemma 2 shows that \(e_i e_j = 0\) . Finally \(e = \sum_{i\in I}e_i\) is idempotent and \(e_i\in \mathbf{A}e\) for all \(i\in I\) , whence \(\mathrm{V}(\mathbf{A}e)\subset \mathbf{Z}_i\) for all \(i\) ; it follows that

\[
\mathbf {V} (\mathbf {A} e) = \varnothing = \mathbf {V} (\mathbf {A}. 1)
\]

and Lemma 2 shows also that \(e = 1\) .

COROLLARY 1. Let A be a ring, r a nil ideal of A and h: A → A/r the canonical homomorphism. For every finite orthogonal family \((e_{i}^{\prime})_{i\in I}\) of idempotents of A/r such that \(\sum_{i\in I} e_{i}^{\prime}=1\) , there exists a finite orthogonal family \((e_{i})_{i\in I}\) of idempotents of A such that \(\sum_{i\in I} e_{i}=1\) and \(h(e_{i})=e_{i}^{\prime}\) for all \(i\in I\) .

We write \(\mathbf{A}' = \mathbf{A} / \mathfrak{r}\) . We know (Remark following Proposition 13) that

\[
^ {a} h \colon \operatorname{Spec} (\mathrm{A} ^ {\prime}) \rightarrow \operatorname{Spec} (\mathrm{A})
\]

is a homeomorphism, every prime ideal of A containing \(\mathfrak{r}\) by hypothesis. Proposition 15 shows that there exists in A a finite orthogonal family \((e_i)_{i\in I}\) of idempotents such that \(\sum_{i\in I}e_i = 1\) and that the image under \(^ah\) of \(\mathrm{V(A'(1 - e_i'))}\) is \(\mathrm{V(A(1 - e_i))}\) . But clearly \(\mathrm{V(A(1 - e_i))}\) is also the image under \(^ah\) of

\[
\mathbf {V} (\mathbf {A} ^ {\prime} (1 - h (e _ {i})))  ;
\]

as \(1 - e_i'\) and \(1 - h(e_i)\) are idempotent, Lemma 2 shows that \(e_i' = h(e_i)\) , whence the corollary.

COROLLARY 2. For the prime spectrum \(\mathbf{X} = \operatorname{Spec}(\mathbf{A})\) of a ring \(\mathbf{A}\) to be connected, it is necessary and sufficient that \(\mathbf{A}\) contain no idempotents other than 0 and 1.

To say that X is not connected means that there exists in X a set which is open and closed and distinct from \(\varnothing\) and X.

